% Generated mechanically from frozen input p12/ja/build/ch112/guide_full.tex.
% Do not edit this generated file; edit the bound locale source instead.

% OK, start here.
%

\setcounter{chapter}{111}
\chapter{文献案内}


\phantomsection
\label{guide-section-phantom}


\section{短い入門論文}
\label{guide-section-short-introductions}


\begin{itemize}
\item Barbara Fantechi: \emph{Stacks for Everybody} \cite{fantechi_stacks}
\item Dan Edidin: \emph{What is a stack?} \cite{edidin_whatis}
\item Dan Edidin: \emph{Notes on the construction of the moduli space of
curves} \cite{edidin_notes}
\item Angelo Vistoli: \emph{Intersection theory on algebraic
stacks and on their moduli spaces}，とりわけその付録．
\cite{vistoli_intersection}
\end{itemize}


\section{古典的参考文献}
\label{guide-section-classics}

\begin{itemize}
\item Mumford: \emph{Picard groups of moduli problems}
\cite{mumford_picard}
\begin{quote}
Mumford はここで「stack」という語を一度も用いていないが，その概念は論文に
暗黙に現れている；彼は楕円曲線のモジュライ・スタックの Picard 群を計算する．
\end{quote}
\item Deligne, Mumford: \emph{The irreducibility of the space of curves of
given genus} \cite{DM}
\begin{quote}
この影響力の大きい論文は，現在では普遍的に Deligne--Mumford スタックと
呼ばれる意味での「代数スタック」，すなわち表現可能な対角射をもち，
スキームによる \'etale 表示を許すスタックを導入する．
多くの基礎的結果が \emph{証明なしに} 述べられている．
この論文ではスタックを用いて，種数 $g$ の曲線のモジュライ空間の
既約性を二通りに証明する．
\end{quote}
\item Artin: \emph{Versal deformations and algebraic stacks}
\cite{ArtinVersal}
\begin{quote}
この論文は，Deligne--Mumford スタックを一般化する「代数スタック」を
導入する．現在これらは通常 \emph{Artin スタック}，すなわち表現可能な
対角射をもち，スキームによる滑らかな表示を許すスタックと呼ばれる．
また，Artin の判定法として知られる変形理論的判定法を与える．
これにより，表示を明示的に構成することなく，与えられたモジュライ・
スタックが Artin スタックであることを証明できる．
\end{quote}
\end{itemize}




\section{書籍とオンライン講義ノート}
\label{guide-section-books}


\begin{itemize}
\item Laumon, Moret-Bailly: \emph{Champs Alg\'ebriques} \cite{LM-B}
\begin{quote}
この書籍は現在，スタックに関する最も網羅的な参考文献であり，多くの
基礎的結果を含む．読者が代数空間に習熟していることを仮定し，
Knutson の書籍 \cite{Kn} を頻繁に参照する．第 12 章には，代数スタックの
lisse-\'etale サイトの関手性に関する誤りがある．しかし Martin Olsson が
その誤りを修補しており（\cite{olsson_sheaves} を参照），残りの章の結果は
（場合によってはわずかな修正を施せば）正しいので，これを心配する必要はない．
\end{quote}
\item The Stacks Project Authors: \emph{Stacks Project}
\cite{stacks-project}.
\begin{quote}
いま読んでいるのがそれである！
\end{quote}
\item Anton Geraschenko:
\emph{Lecture notes for Martin Olsson's class on stacks} \cite{olsson_stacks}
\begin{quote}
この講義では，代数スタックを導入する前に代数空間の理論を体系的に展開する
（代数スタックが初めて定義されるのは講義 27 である！）．基本的性質に加えて，
Deligne--Mumford であることと不分岐な対角射をもつこととの同値性，
Artin スタック上の lisse-\'etale サイト，準連接層の理論，Keel--Mori 定理，
コホモロジー的降下，およびジェルブ（と Brauer 群との関係）を扱う．
演習問題もいくつか含まれる．
\end{quote}
\item
Behrend, Conrad, Edidin, Fantechi, Fulton, G\"ottsche, and Kresch:
\emph{Algebraic stacks}，現在執筆中の書籍のオンライン講義ノート
\cite{stacks_book}
\begin{quote}
この書籍の目的は，高度な予備知識を仮定せず，例と応用に重点を置いて
スタックへの親しみやすい入門を与えることである．\cite{LM-B} と異なり，
読者が代数空間の理論をすでに消化しているとは仮定しない．
その代わり，代数空間を特別な場合として Deligne--Mumford スタックを導入し，
\cite{DM} の主張を証明するのに十分な理論を展開することを目標の一部とする．
Artin スタックの一般理論は第 2 部で展開される予定である．
現在 Kresch のウェブサイトで入手できるのは書籍の一部分だけである．
\end{quote}
\item Olsson, Martin: \emph{Algebraic spaces and stacks}, \cite{olsson_book}
\begin{quote}
Hartshorne の代数幾何学の書籍を修得した程度から始められる，
代数空間と代数スタックへの非常に推奨できる入門書である．
\end{quote}
\end{itemize}


\section{スタックの基礎に関する関連文献}
\label{guide-section-related}


\begin{itemize}
\item Vistoli:
\emph{Notes on Grothendieck topologies, fibered categories and descent theory}
\cite{vistoli_fga}
\begin{quote}
fpqc 位相におけるファイバー化圏，スタック，降下理論についての有用な事実と，
厳密な証明を含む．
\end{quote}
\item Knutson: \emph{Algebraic Spaces} \cite{Kn}
\begin{quote}
Michael Artin の下で書かれた博士論文から発展したこの書籍は，
代数空間の理論の基礎を含む．\cite{LM-B} はこの書籍を頻繁に参照する．
代数空間に関する Artin の論文
\cite{Artin-Algebraic-Approximation},
\cite{ArtinI}, \cite{Artin-Implicit-Function},
\cite{ArtinII}, \cite{Artin-Construction-Techniques},
\cite{Artin-Algebraic-Spaces}, \cite{Artin-Theorem-Representability}, および
\cite{ArtinVersal}
も参照されたい．
\end{quote}
\item Grothendieck et al, \emph{Th\'eorie des Topos et Cohomologie \'Etale des
Sch\'emas I, II, III}，SGA4 としても知られる \cite{SGA4}
\begin{quote}
第 1 巻は，宇宙，サイト，ファイバー化圏に関する多くの一般的事実を含む．
「champ」（フランス語で「stack」）という語は，Deligne の Expos\'e XVIII
に現れる．
\end{quote}
\item Jean Giraud: \emph{Cohomologie non ab\'elienne} \cite{giraud}
\begin{quote}
この書籍は，一般のサイト上のファイバー化圏，スタック，トーサー，
ジェルブを論じるが，代数スタックは論じない．たとえば $G$ が $X$ 上の
可換群の層ならば，$H^1(X, G)$ を $G$-トーサーと同一視できるのと同様に，
$H^2(X, G)$ を適切に定義された $G$-ジェルブの集合と同一視できる．
$G$ が可換でないときには，$H^2(X, G)$ は $G$-ジェルブの集合として定義される．
\end{quote}
\item Kelly and Street: \emph{Review of the elements of 2-categories}
\cite{kelly-street}
\begin{quote}
% P12-ERR-0556: frozen English has compound agreement/syntax defects here; see ../../control/ERRATA_CANDIDATES.jsonl.
スタックの圏は 2-圏をなすが，すべての 2-射が可逆であるという単純な型の
2-圏である．これは一般の 2-圏に関する参考文献である．著者はこれを
使ったことがないので，どの程度有用かを述べることはできない．
なお，\cite{stacks-project} も 2-圏の基礎をいくらか含む．
\end{quote}
\end{itemize}

\section{文献中の論文}
\label{guide-section-papers}

\noindent
以下は，スタックおよび代数空間に関する基礎的結果を含む研究論文の一覧である．
要約の意図は，各論文の結果のうちスタック理論に寄与するものだけを示すことに
ある；多くの場合，これらの結果は論文の主目的に対して副次的である．
論文をいくつかの範疇に分けるが，複数の範疇に属するものもある．


\subsection{変形理論と代数スタック}
\label{guide-subsection-deformation-theory}

\noindent
Artin による最初の 3 論文はスタックについて何も含まないが，強力な結果を
含んでおり，最初の 2 論文は \cite{ArtinVersal} に不可欠である．
\begin{itemize}
\item Artin: \emph{Algebraic approximation of structures over
complete local rings} \cite{Artin-Algebraic-Approximation}
\begin{quote}
穏やかな仮定の下で，任意の有効形式変形を近似できることが証明される：
$F: (\Sch/S) \to (\textit{Sets})$ を局所有限表示な反変関手とし，
$S$ は体上有限型または優秀な DVR，$s \in S$ とし，
$\hat{\xi} \in F(\hat{\mathcal{O}}_{S, s})$ を有効形式変形とする．
このとき任意の $n > 0$ に対し，剰余体上自明な \'etale 近傍
$(S', s') \to (S, s)$ と $\xi' \in F(S')$ が存在し，$\xi'$ と
$\hat{\xi}$ は $n$ 次まで一致する（すなわち
$F(\mathcal{O}_{S, s} / \mathfrak m^n)$ における制限が同じである）．
\end{quote}
\item
Artin: \emph{Algebraization of formal moduli I} \cite{ArtinI}
\begin{quote}
穏やかな仮定の下で，任意の有効形式半普遍変形が代数化可能であることが
証明される．$F: (\Sch/S) \to (\textit{Sets})$ を局所有限表示な反変関手とし，
$S$ は体上有限型または優秀な DVR，$s \in S$ は局所閉点，
$\hat A$ は剰余体 $k'$ が $k(s)$ の有限拡大である完備 Noether
局所 $\mathcal{O}_S$-代数とし，$\hat{\xi} \in F(\hat A)$ を
$\xi_0 \in F(k')$ の有効形式半普遍変形とする．このとき $S$ 上有限型な
スキーム $X$，剰余体 $k(x) = k'$ をもつ閉点 $x \in X$，
および元 $\xi \in F(X)$ が存在し，同型
$\hat{\mathcal{O}}_{X, x} \cong \hat{A}$ があって，各
$F(\hat A / \mathfrak m^n)$ における $\xi$ と $\hat{\xi}$ の制限を
同一視する．$\hat{\xi}$ が普遍変形ならば，この代数化は一意である．
Hilbert スキームと Picard スキームの表現可能性への応用が与えられる．
\end{quote}
\item Artin: \emph{Algebraization of formal moduli. II}
\cite{ArtinII}
\begin{quote}
% P12-ERR-0557: frozen English switches from Y' to Y in the contracted subset; see ../../control/ERRATA_CANDIDATES.jsonl.
大まかにいえば，閉部分集合 $Y' \subset X'$ を $Y'$ の形式的局所近傍で
収縮できるならば，$Y$ を収縮する大域射 $X' \to X$ が存在し，
$X$ は代数空間となることが示される．
\end{quote}
\item
Artin: \emph{Versal deformations and algebraic stacks} \cite{ArtinVersal}
\begin{quote}
この画期的な論文は，\cite{Artin-Algebraic-Approximation} と
\cite{ArtinI} における Artin の研究に基づく．スタックが代数的であることを，
変形理論的性質の検証によって証明できる Artin の判定法を導入する．
より正確には（とはいえ，それほど正確ではないが），Artin は極限を保存する
スタック $\mathcal{X}$ の点 $x \in \mathcal{X}(k)$ の近傍における表示を
次のように構成する．スタック $\mathcal{X}$ が Schlessinger の判定法
（\cite{Sch}）を満たすと仮定すると，$x$ の形式的半普遍変形
$\hat{\xi} \in \lim \mathcal{X}(\hat A / \mathfrak m^n)$ が存在する．
形式変形が有効である，すなわち
$\mathcal{X}(\hat{A}) \to \lim \mathcal{X}(\hat A / \mathfrak m^n)$
が全単射であると仮定すれば，有効形式半普遍変形
$\xi \in \mathcal{X}(\hat A)$ が得られる．\cite{ArtinI} の結果を用いて，
有限型スキーム $U$ と元 $\xi_U: U \to \mathcal{X}$ を構成でき，これは
$x$ 上の点 $u \in U$ において形式的半普遍である．さらに，
$\mathcal{X}$ が一定の条件，すなわち \'etale 局所化および完備化との
両立性と構成可能性条件を満たす変形・障害理論をもつと仮定する．
すると第 4 節で，形式的半普遍性が開条件であることが示されるので，
$U$ を縮小した後に $U \to \mathcal{X}$ は滑らかとなる．
Artin はさらに，スキームによる fppf 表示を許す任意のスタックが
スキームによる滑らかな表示を許すことも証明する．したがって特に，
平坦，分離的，有限表示な群スキームによる商スタックを形成できる．
\end{quote}
\item Conrad, de Jong: \emph{Approximation of Versal Deformations}
\cite{conrad-dejong}
\begin{quote}
この論文は Popescu の強力な結果を適用して Artin の代数化結果に至る
一つの方法を与える：$A$ が Noether 環，$B$ が Noether
$A$-代数ならば，写像 $A \to B$ が正則射であることと，$B$ が滑らかな
$A$-代数の帰納極限であることは同値である．Popescu の結果から，
任意の優秀スキーム上の Artin 近似が従うことは難しくない
（優秀性の仮定により，局所環 $A$ に対して Hensel 化から完備化への写像
$A^{\text{h}} \to \hat A$ は正則となる）．この論文は Popescu の結果を用い，
\cite{ArtinI} の主定理を「亜群」へ一般化する．この一般化は任意の
優秀基礎スキーム上，任意の点 $s \in S$ に対して有効である．
特に \cite{ArtinVersal} の結果は任意の優秀基礎上で成り立つ．
また，代数化の \'etale 局所的一意性と，対象の自己同型群が代数化の
Hensel 化に自然に作用するかどうかを論じる．
\end{quote}
\item Jason Starr: \emph{Artin's axioms, composition, and moduli spaces}
\cite{starr_artin}
\begin{quote}
この論文は，代数化に関する Artin の公理が 1-射の合成と両立することを示す．
\end{quote}
\item Martin Olsson: \emph{Deformation theory of representable
morphism of algebraic stacks} \cite{olsson_deformation}
\begin{quote}
% P12-ERR-0558: frozen English title/body have singular-agreement issues; see ../../control/ERRATA_CANDIDATES.jsonl.
スキームの射に対する標準的な変形理論の結果を，余接複体を用いて
代数スタックの表現可能射へ一般化する．これらの結果を Illusie の一般理論の
帰結とみなすことはできない．表現可能射 $X \to \mathcal{X}$ の余接複体は，
環付きトポスの射の余接複体によっては定義されないからである
（lisse-\'etale サイトには関手性がない）．
\end{quote}
\end{itemize}

\subsection{粗モジュライ空間}
\label{guide-subsection-coarse-moduli-spaces}

\noindent
粗モジュライ空間を論じる論文．
\begin{itemize}
\item Keel, Mori: \emph{Quotients in Groupoids} \cite{K-M}
\begin{quote}
分離的 Deligne--Mumford スタックが粗モジュライ空間をもつことは，
長らく「民間伝承」だったらしい．ここでは次の定理の厳密な
（ただし簡潔な）証明が与えられる：$\mathcal{X}$ を Noether 基礎スキーム
上局所有限型な Artin スタックとし，慣性スタック
$I_\mathcal{X} \to \mathcal{X}$ が有限であるとする．このとき，
$\phi$ が分離的で，$Y$ が $S$ 上局所有限型な代数空間であるような
粗モジュライ空間 $\phi : \mathcal{X} \to Y$ が存在する．
慣性が有限という仮定はまさに適切な条件である：$\phi$ が分離的な
粗モジュライ空間 $\phi : \mathcal{X} \to Y$ が存在することと，
慣性が有限であることは同値である．
\end{quote}
\item Conrad: \emph{The Keel-Mori Theorem via Stacks} \cite{conrad}
\begin{quote}
Keel と Mori の論文 \cite{K-M} は亜群の言葉で書かれており，
理解するのが難しいと感じる人もいる．Brian Conrad は，高度な
スタックの言葉を用いるものの，かなり明快な証明のスタック理論版を与える．
Conrad はさらに Noether 性の仮定を除く．
\end{quote}
\item Rydh: \emph{Existence of quotients by finite groups and coarse moduli
spaces} \cite{rydh_quotients}
\begin{quote}
Rydh は，$\mathcal{X}$ がある基礎上有限表示であるという
\cite{K-M} と \cite{conrad} の仮定を除く．
\end{quote}
\item
Abramovich, Olsson, Vistoli: \emph{Tame stacks in positive characteristic}
\cite{tame}
\begin{quote}
彼らは \emph{tame Artin スタック}を，有限な慣性をもつ Artin スタックで，
粗モジュライ空間を $\phi : \mathcal{X} \to Y$ とするとき
$\phi_*$ が準連接層上で完全となるものとして定義する．
有限な慣性をもつ Artin スタックについて，次が同値であることを証明する：
$\mathcal{X}$ が tame であること；$\mathcal{X}$ の安定化群が線形簡約的で
あること；粗モジュライ空間上 \'etale 局所的に，$\mathcal{X}$ が
線形簡約群スキームによるアフィン・スキームの商であること．
tame Artin スタックの粗モジュライ空間は特によい性質をもつ．
たとえば，粗モジュライ空間の形成は任意の基底変換と可換するが，
有限な慣性をもつ一般の Artin スタックの粗モジュライ空間は，
平坦な基底変換とのみ可換する．
\end{quote}
\item Alper: \emph{Good moduli spaces for Artin stacks} \cite{alper_good}
\begin{quote}
無限なアフィン安定化群をもつ一般の Artin スタック
（これは必然的に非分離的である）には，粗モジュライ空間が存在しないことが
多い．最も簡単な例は $[\mathbf{A}^1 / \mathbf{G}_m]$ である．
ここでは，準コンパクト射 $\phi : \mathcal{X} \to Y$ が
\emph{良モジュライ空間}であるとは，
$\mathcal{O}_Y \to \phi_* \mathcal{O}_\mathcal{X}$ が同型で，
$\phi_*$ が準連接層上で完全であることと定義する．
この概念は \cite{tame} の tame Artin スタックを一般化すると同時に，
Mumford の幾何学的不変式論を包含する：簡約群 $G$ が
$X \subset \mathbf{P}^n$ に線形に作用するとき，半安定部分の商スタックから
GIT 商への射 $[X^{ss}/G] \to X//G$ は良モジュライ空間である．
良モジュライ空間は多くのよい幾何学的性質をもつ：
(1) $\phi$ は全射，普遍閉，普遍沈降的である；
(2) $\phi$ は $Y$ の点と $\mathcal{X}$ の点を閉包同値まで同一視する；
(3) $\phi$ は代数空間への写像に関して普遍的である；
(4) 良モジュライ空間は任意の基底変換で安定である；
(5) Artin スタック上のベクトル束が良モジュライ空間へ降下するための
必要十分条件は，閉点における表現が自明であることである．
\end{quote}
\end{itemize}


\subsection{交叉理論}
\label{guide-subsection-intersection-theory}

\noindent
代数スタック上の交叉理論を論じる論文．
\begin{itemize}
\item
Vistoli: \emph{Intersection theory on algebraic stacks and on their moduli
spaces} \cite{vistoli_intersection}
\begin{quote}
この論文は，Deligne--Mumford スタックに対する有理係数交叉理論の基礎を
展開する．$\mathcal{X}$ が分離的 Deligne--Mumford スタックならば，
有理係数の Chow 群 $\CH_*(\mathcal{X})$ は，次元 $k$ の整閉部分スタックが
生成する自由アーベル群を有理同値で割ったものとして定義される．
% P12-ERR-0559: frozen English mixes CH_* and dimension k, then maps to CH_k; see ../../control/ERRATA_CANDIDATES.jsonl.
平坦引き戻し，固有順像，および正則局所埋め込みに対する一般化 Gysin
準同型が存在する．$\phi : \mathcal{X} \to Y$ がモジュライ空間
（すなわち，幾何学的点上全単射な固有射）ならば，誘導される順像
$\CH_*(\mathcal{X}) \to \CH_k(Y)$ は同型である．
\end{quote}
\item Edidin, Graham: \emph{Equivariant Intersection Theory}
\cite{edidin-graham}
\begin{quote}
この論文の目的は，代数群 $G$ の代数空間 $X$ への作用の商スタック
$[X/G]$ に対する整数係数交叉理論，言い換えれば $X$ の
$G$-同変交叉理論を展開することである．不変サイクルだけを用いて定義した
同変 Chow 群からは，よい性質をもつ理論は得られない．
そこで $BG$ の場合の Totaro の定義を一般化し，ベクトル束
$V \to X$ に対して自然な同型 $\CH_i(X) \cong \CH_i(V)$ があることに
動機づけられて，著者らは $\CH_i^G(X)$ を次のように定義する．
$\dim(X) = n$, $\dim(G) = g$ とする．各 $i$ に対し，
$l$ 次元 $G$-表現 $V$ を選び，$G$ が開部分集合 $U \subset V$ 上自由に
作用し，その補集合の余次元が $d > n - i$ であるようにする．すると
$X_G = [X \times U / G]$ は代数空間となる（スキームとなるように選ぶこと
さえできる）．このとき
$\CH_i^G(X) = \CH_{i + l - g}(X_G)$ と定義する．商スタックについては
$\CH_i( [X/G]) = \CH_{i + g}^G(X) = \CH_{i + l}(X_G)$ と定義する．
特に $i > \dim [X/G] = n - g$ ならば $\CH_i([X/G]) = 0$ であるが，
$i \ll 0$ に対して非零となることがある．たとえば
$i \le 0$ に対して $\CH_i(B \mathbf{G}_m) = \mathbf{Z}$ である．
これらの同変 Chow 群が通常の Chow 群と同じ関手的性質をもつことを示す．
さらに，$[X / G] \cong [Y / H]$ ならば
$\CH_i([X/G]) = \CH_i([Y/H])$ であることを示すので，定義はスタックを
商スタックとしてどのように表示するかに依存しない．
\end{quote}
\item
Kresch: \emph{Cycle Groups for Artin Stacks} \cite{kresch_cycle}
\begin{quote}
Kresch は任意の Artin スタックに対して Chow 群を定義する．
商スタックの場合には \cite{edidin-graham} における Edidin と Graham の
定義と一致する．アフィン安定化群をもつ代数スタックでは，
この理論は通常の性質を満たす．
\end{quote}
\item Behrend and Fantechi: \emph{The intrinsic normal cone}
\cite{behrend-fantechi}
\begin{quote}
Behrend と Fantechi は，Li と Tian による構成を一般化して，
Deligne--Mumford スタックに対する仮想基本類を構成する．
\end{quote}
\end{itemize}

\subsection{商スタック}
\label{guide-subsection-quotient-stacks}

\noindent
商スタック\footnote{文献では \emph{quotient stack} が，任意の平坦群
スキームによるものではなく，$X$ が代数空間で $G$ が
$\text{GL}_n$ の部分群スキームである形 $[X/G]$ のスタックを
意味することが多い．}は Artin スタックの非常に重要な部分類をなし，
代数幾何学者が研究するほとんどすべてのモジュライ・スタックを含む．
商スタック $[X/G]$ の幾何学は $X$ の $G$-同変幾何学である．
商スタックに対して性質を示す方が容易なことが多く，商スタックについて
のみ知られている結果もある．以下の論文は次の問題を扱う：
代数スタックはいつ大域的商スタックとなるか？
代数スタックは「局所的に」商スタックとなるか？
\begin{itemize}
\item Laumon, Moret-Bailly: \cite[Chapter 6]{LM-B}
\begin{quote}
第 6 章は代数スタックの局所構造および大域構造に関するいくつかの事実を
含む．$S$ 上の代数スタック $\mathcal{X}$ が，$Y$ が代数空間
（それぞれスキーム，アフィン・スキーム）で $G$ が有限群であるような
商スタック $[Y/G]$ となるための必要十分条件は，代数空間
（それぞれスキーム，アフィン・スキーム）$Y'$ と有限 \'etale 射
$Y' \to \mathcal{X}$ が存在することであると証明される．
また，$S$ 上の任意の Deligne--Mumford スタックと
$x : \Spec(K) \to \mathcal{X}$ に対し，$S$ 上のアフィン・スキームに
作用する有限群 $G$ と，表現可能，\'etale，分離的な射
$\phi : [X/G] \to \mathcal{X}$ が存在して，
$\Spec(K) = [X/G] \times_\mathcal{X} \Spec(K)$ となることが示される．
幾何学的に連結なファイバーをもつ表示の存在も詳しく論じられる．
\end{quote}
\item
Edidin, Hassett, Kresch, Vistoli:
\emph{Brauer Groups and Quotient stacks} \cite{ehkv}
\begin{quote}
まず，与えられた代数スタック（この論文では常に Noether スキーム上
有限型と仮定する）がいつ商スタックとなるかに関する，いくつかの基礎的な
（ただしそれほど難しくない）事実を示す．代数スタック $\mathcal{X}$ について，
次は同値である：$\mathcal{X}$ が商スタックである；ベクトル束
$V \to \mathcal{X}$ が存在し，各幾何学的点で安定化群がファイバーに忠実に
作用する；ベクトル束 $V \to \mathcal{X}$ と局所閉部分スタック
$V^0 \subset V$ が存在して，$V^0$ が表現可能かつ $\mathcal{X}$ 上全射である．
代数空間による有限平坦被覆が存在すれば，代数スタックが商スタックであることも
示す．一般点で安定化群が自明な滑らかな Deligne--Mumford スタックはすべて
商スタックである．Noether スキーム $X$ 上の $\mathbf{G}_m$-ジェルブで，
$\beta \in H^2(X, \mathbf{G}_m)$ に対応するものが商スタックであるための
必要十分条件は，$\beta$ が Brauer 写像
$\text{Br}(X) \to \text{Br}'(X)$ の像に属することである．
これを用いて，商スタックでない非分離的 Deligne--Mumford スタックを構成する．
\end{quote}
\item Totaro: \emph{The resolution property for schemes and stacks}
\cite{totaro_resolution}
\begin{quote}
すべての連接層がベクトル束の商であるとき，スタックは
resolution property（解消性質）をもつという．
第 1 の主定理は，$\mathcal{X}$ が閉点でアフィン安定化群をもつ正規
Noether 代数スタックならば，次が同値であるというものである：
(1) $\mathcal{X}$ は解消性質をもつ；
(2) $Y$ が準アフィンであるように
$\mathcal{X} = [Y/\text{GL}_n]$ と表せる．
$\mathcal{X}$ が体上有限型の場合，(1), (2) はさらに次と同値である：
(3) $G$ が $k$ 上有限型なアフィン群スキームであるように
$\mathcal{X} = [\Spec(A)/G]$ と表せる．商スタックが解消性質をもつことは
Thomason により証明された．第 2 の主定理は次の通りである：
$\mathcal{X}$ を体上の滑らかな Deligne--Mumford スタックとし，
有限かつ一般点で自明な安定化群 $I_\mathcal{X} \to \mathcal{X}$ をもち，
その粗モジュライ空間がアフィン対角射をもつスキームであるとする．
このとき $\mathcal{X}$ は解消性質をもつ．もう一つの興味深い結果は，
$\mathcal{X}$ が解消性質を満たす Noether 代数スタックならば，
$\mathcal{X}$ がアフィン対角射をもつことと，その閉点がアフィン安定化群を
もつことが同値である，というものである．
\end{quote}
\item Kresch: \emph{On the Geometry of Deligne-Mumford Stacks}
\cite{kresch_geometry}
\begin{quote}
この論文は（体上有限型な）Deligne--Mumford スタックの一般構造定理を
要約し，商スタックに関する興味深い結果をいくつか含む．
準射影的粗モジュライ空間をもつ，滑らかで分離的かつ一般点で tame な
Deligne--Mumford スタックはすべて，$Y$ が準射影的で $G$ が代数群である
商スタック $[Y/G]$ となることが示される．Deligne--Mumford スタック
$\mathcal{X}$ の粗モジュライ空間がスキームであるとする．
このとき $\mathcal{X}$ が Zariski 局所的に商スタックで
あるための必要十分条件は，有限群によるスキームのスタック商からなる
Zariski 開被覆を許すことである．標数 0 の体上固有な
Deligne--Mumford スタック $\mathcal{X}$ の粗モジュライ空間を $Y$ とする．
このとき次は同値である：$Y$ は射影的で $\mathcal{X}$ は商スタックである；
$Y$ は射影的で $\mathcal{X}$ は生成層をもつ；
$\mathcal{X}$ は射影的粗モジュライ空間をもつ滑らかで固有な
Deligne--Mumford スタックへの閉埋め込みを許す．
これを動機として，Deligne--Mumford スタックが \emph{射影的}であるとは，
射影的粗モジュライ空間をもつ滑らかで固有な Deligne--Mumford スタックへの
閉埋め込みが存在することと定義する．
\end{quote}
\item Kresch, Vistoli \emph{On coverings of Deligne-Mumford stacks and
surjectivity of the Brauer map} \cite{kresch-vistoli}
\begin{quote}
標数 0 かつ固定した $n$ に対し，次の二つの主張が同値であることが示される：
(1) 次元 $n$ の滑らかな Deligne--Mumford スタックはすべて商スタックである；
(2) 次元 $n$ の滑らかなスキームに対し，Azumaya Brauer 群と
コホモロジー的 Brauer 群が一致する．
\end{quote}
\item Kresch: \emph{Cycle Groups for Artin Stacks} \cite{kresch_cycle}
\begin{quote}
体上有限型でアフィン安定化群をもつ被約 Artin スタックが，
商スタックによる層化を許すことが示される．
\end{quote}
\item Abramovich-Vistoli:
\emph{Compactifying the space of stable maps} \cite{abramovich-vistoli}
\begin{quote}
% P12-ERR-0560: frozen English omits a predicate/article in the Lemma 2.2.3 sentence; see ../../control/ERRATA_CANDIDATES.jsonl.
補題 2.2.3 は，任意の分離的 Deligne--Mumford スタックが，粗モジュライ空間上
\'etale 局所的に，$U$ がアフィンで $G$ が有限群である商スタック
$[U/G]$ となることを示す．\cite[Theorem 2.12]{olsson_homstacks} は，
この議論において $G$ が実際には安定化群であることを示す．
\end{quote}
\item Abramovich, Olsson, Vistoli:
\emph{Tame stacks in positive characteristic} \cite{tame}
\begin{quote}
% P12-ERR-0561: frozen English says “quotient stack of an affine” with the noun omitted; see ../../control/ERRATA_CANDIDATES.jsonl.
この論文は，tame Artin スタックが粗モジュライ空間上 \'etale 局所的に，
安定化群によるアフィンの商スタックであることを示す．
\end{quote}
\item Alper: \emph{On the local quotient structure of Artin stacks}
\cite{alper_quotient}
\begin{quote}
Artin スタック $\mathcal{X}$ と線形簡約的な安定化群をもつ閉点
$x \in \mathcal{X}$ に対し，$V$ が代数空間であるような \'etale 射
$[V/G_x] \to \mathcal{X}$ が存在する，という予想が述べられる．
この予想を支持するいくつかの証拠が与えられる．
\cite{tame} の着想に基づく簡単な変形理論の議論により，
形式的局所には正しいことが示される．Luna の \'etale slice 定理の
スタック理論的証明も与えられ，$G$ が線形簡約的で
$\mathcal{X} = [\Spec(A)/G]$ であるスタックに対して，
GIT 商 $\Spec(A^G)$ 上 \'etale 局所的に，$\mathcal{X}$ は
安定化群による商スタックであることが証明される．
\end{quote}
\end{itemize}

\subsection{コホモロジー}
\label{guide-subsection-cohomology}

\noindent
代数スタック上の層のコホモロジーを論じる論文．
\begin{itemize}
\item Olsson: \emph{Sheaves on Artin stacks} \cite{olsson_sheaves}
\begin{quote}
この論文は準連接層および構成可能層の理論を展開し，基本的な
コホモロジー的性質を証明する．また \cite{LM-B} にある
lisse-\'etale サイトの関手性に関する誤りを訂正する．
余接複体が構成される．さらに，次の定理が証明される：
固有射に対する Grothendieck の基本定理，Grothendieck の存在定理，
Zariski の連結性定理，ならびに連接層および構成可能層の固有順像に対する
有限性定理．
\end{quote}
\item Behrend: \emph{Derived $l$-adic categories for algebraic stacks}
\cite{behrend_derived}
\begin{quote}
代数スタックに対する Lefschetz 跡公式を証明する．
\end{quote}
\item Behrend: \emph{Cohomology of stacks} \cite{behrend_cohomology}
\begin{quote}
可微分スタックに対する de Rham コホモロジーと，位相スタックに対する
特異コホモロジーを定義する．
\end{quote}
\item Faltings:
\emph{Finiteness of coherent cohomology for proper fppf stacks}
\cite{faltings_finiteness}
\begin{quote}
固有射による連接層の順像の連接性を証明する．
\end{quote}
\item Abramovich, Corti, Vistoli:
\emph{Twisted bundles and admissible covers} \cite{acv}
\begin{quote}
付録は tame Deligne--Mumford スタックの \'etale コホモロジーに対する
固有基底変換定理を含む．
\end{quote}
\end{itemize}


\subsection{スキームによる有限被覆の存在}
\label{guide-subsection-finite-covers}

\noindent
Deligne--Mumford スタックのスキームによる有限被覆の存在は重要な結果である．
Deligne--Mumford スタック上の交叉理論では，表現可能でない射に対する
固有順像を定義するための本質的な材料となる．
\emph{滑らかな}スキームによる有限被覆の存在に依存する
$\overline{\mathcal{M}}_g$ に関する結果がいくつかあり，その存在は
Looijenga によって証明された．この方向で最初の結果は，同変の場合を扱う
\cite[Theorem 6.1]{seshadri_quotients} かもしれない．
\begin{itemize}
\item Vistoli: \emph{Intersection theory on algebraic stacks and on their
moduli spaces} \cite{vistoli_intersection}
\begin{quote}
$\mathcal{X}$ がモジュライ空間をもつ Deligne--Mumford スタック
（すなわち，幾何学的点上全単射な固有射をもつもの）ならば，
スキーム $X$ からの有限射 $X \to \mathcal{X}$ が存在する．
\end{quote}
\item Laumon, Moret-Bailly: \cite[Chapter 16]{LM-B}
\begin{quote}
Zariski の主定理の応用として，定理 16.6 は次を示す：
$\mathcal{X}$ が Noether スキーム上有限型な Deligne--Mumford スタック
ならば，$Z$ がスキームであるような有限，全射，一般点で \'etale な射
$Z \to \mathcal{X}$ が存在する．系 16.6.2 では，任意の Noether 正規
代数空間が，正規スキーム $X$ に作用する有限群 $G$ による代数空間商
$X'/G$ と同型であることも示される．
\end{quote}
\item Edidin, Hassett, Kresch, Vistoli:
\emph{Brauer Groups and Quotient stacks} \cite{ehkv}
\begin{quote}
定理 2.7 は次を述べる：$\mathcal{X}$ を Noether 基礎スキーム $S$ 上
有限型な代数スタックとする．このとき対角射
$\mathcal{X} \to \mathcal{X} \times_S \mathcal{X}$ が準有限であることと，
スキーム $X$ からの有限全射
% P12-ERR-0562: frozen English target F conflicts with the surrounding Xcal; see ../../control/ERRATA_CANDIDATES.jsonl.
$X \to F$ が存在することは同値である．
\end{quote}
\item Kresch, Vistoli: \emph{On coverings of Deligne-Mumford stacks and
surjectivity of the Brauer map} \cite{kresch-vistoli}
\begin{quote}
体上有限型で準射影的粗モジュライ空間をもつ滑らかな分離的
Deligne--Mumford スタックはすべて，滑らかな準射影スキームによる
有限平坦被覆を許すことが証明される．
\end{quote}
\item Olsson: \emph{On proper coverings of Artin stacks} \cite{olsson_proper}
\begin{quote}
$\mathcal{X}$ が $S$ 上分離的かつ有限型な Artin スタックならば，
$S$ 上準射影的なスキーム $X$ からの固有全射
$X \to \mathcal{X}$ が存在することを証明する．応用として Olsson は，
固有射による層の順像の連接性および構成可能性を証明する．
さらに，固有 Artin スタックに対する Grothendieck の存在定理を証明する．
\end{quote}
\item Rydh: \emph{Noetherian approximation of algebraic spaces and stacks}
\cite{rydh_approx}
\begin{quote}
この論文の定理 B は次の通りである．$X$ を，準有限かつ分離的な対角射をもつ
準コンパクト代数スタック（それぞれ，準コンパクトかつ分離的な対角射をもつ
準コンパクト Deligne--Mumford スタック）とする．このときスキーム $Z$ と，
有限，有限表示，全射な射 $Z \to X$ が存在し，これは稠密な準コンパクト
開部分スタック $U \subset X$ 上で平坦（それぞれ \'etale）である．
\end{quote}
\end{itemize}

\subsection{リジディフィケーション}
\label{guide-subsection-rigidification}

\noindent
リジディフィケーションは，慣性から平坦部分群を取り除く過程である．
たとえば $X$ が射影多様体ならば，Picard スタックから Picard スキームへの
射は，自己同型群 $\mathbf{G}_m$ のリジディフィケーションである．
\begin{itemize}
\item Abramovich, Corti, Vistoli:
\emph{Twisted bundles and admissible covers} \cite{acv}
\begin{quote}
$\mathcal{X}$ を $S$ 上の代数スタック，$H$ を $S$ 上の平坦，有限表示，
分離的な群スキームとする．任意の対象 $\xi \in \mathcal{X}(T)$ に対し，
埋め込み $H(T) \hookrightarrow
\text{Aut}_{\mathcal{X}(T)}(\xi)$ があり，引き戻しと次の意味で両立すると
仮定する：$f: T \rightarrow T'$ 上の任意の射
$\phi : \xi \rightarrow \xi'$ と任意の $g \in H(T')$ に対して，
$g \circ \phi = \phi \circ f^*g$ が成り立つ．
このとき代数スタック $\mathcal{X}/H$ と射
$\rho : \mathcal{X} \rightarrow \mathcal{X}/H$ が存在し，
この射は fppf ジェルブである．さらに各
$\xi \in \mathcal{X}(T)$ に対し，射
$\text{Aut}_{\mathcal{X}(T)} (\xi)
\rightarrow \text{Aut}_{\mathcal{X}/H (T)} (\xi) $
は全射で，核は $H(T)$ である．
\end{quote}
\item Romagny: \emph{Group actions on stacks and applications}
\cite{romagny_actions}
\begin{quote}
群作用がリジディフィケーションに関してどのように振る舞うかを論じる．
\end{quote}
\item Abramovich, Graber, Vistoli:
\emph{Gromov-Witten theory for Deligne-Mumford stacks} \cite{agv}
\begin{quote}
付録は \cite{acv} におけるリジディフィケーションを二つの別の解釈とともに
要約する．この論文は，閉部分スタックに沿って代数スタックを貼り合わせる構成，
および直線束の根を取る構成も含む．
\end{quote}
\item
Abramovich, Olsson, Vistoli: \emph{Tame stacks in positive characteristic}
(\cite{tame})
\begin{quote}
付録は，慣性の平坦部分群スタック $H \subset I_\mathcal{X}$ が正規だが
必ずしも中心的ではない，より複雑な状況を扱う．
\end{quote}
\end{itemize}

\subsection{スタック曲線}
\label{guide-subsection-stacky-curves}

\noindent
スタック曲線を論じる論文．
\begin{itemize}
\item Abramovich, Vistoli: \emph{Compactifying the space of stable maps}
\cite{abramovich-vistoli}
\begin{quote}
この論文は \emph{twisted 曲線}を導入する．安定曲線から代数スタックへの
安定写像のモジュライ空間は，一般にはコンパクトでない．twisted 曲線からの
写像を用いて，標的が射影的粗モジュライ空間をもつ tame
Deligne--Mumford スタックであるとき固有となるモジュライ・スタックを
著者らは構成する．
\end{quote}
\item Behrend, Noohi: \emph{Uniformization of Deligne-Mumford curves}
\cite{behrend-noohi}
\begin{quote}
Deligne--Mumford 解析曲線の一意化定理を証明する．
\end{quote}
\end{itemize}

\subsection{Hilbert, Quot, Hom および branchvariety スタック}
\label{guide-subsection-hilbert-quot-hom}

\noindent
Hilbert スキームなどを論じる論文．
\begin{itemize}
\item Vistoli: \emph{The Hilbert stack and the theory of moduli of families}
\cite{vistoli_hilbert}
\begin{quote}
$\mathcal{X}$ を，局所 Noether かつ局所分離的な代数空間 $S$ 上の，
分離的で局所有限型な代数スタックとする．Vistoli は，固有スキームからの
有限不分岐射をパラメータ化する Hilbert スタック
$\mathcal{H}\text{ilb}(\mathcal{F} / S)$ を定義する．
$\mathcal{H}\text{ilb}(\mathcal{F} / S)$ が代数スタックであることは
証明なしに主張される．その帰結として，上の $\mathcal{X}$ に対し，
$T$ が $S$ 上固有かつ平坦ならば Hom スタック
$\mathcal{H} \text{om}_S(T, \mathcal{X})$ が代数スタックであることが
証明される．
\end{quote}
\item Olsson, Starr: \emph{Quot functors for Deligne-Mumford stacks}
\cite{olsson-starr}
\begin{quote}
$\mathcal{X}$ を代数空間 $S$ 上分離的かつ局所有限表示な
Deligne--Mumford スタックとし，$\mathcal{F}$ を局所有限表示な
$\mathcal{O}_\mathcal{X}$-加群とする．このとき Quot 関手
$\text{Quot}(\mathcal{F} / \mathcal{X} / S)$ は，$S$ 上分離的かつ
局所有限表示な代数空間によって表現される．この論文は生成層も定義し，
有限群によるスキームの大域的商スタックである tame かつ分離的な
Deligne--Mumford スタックに対し，生成層の存在を証明する．
\end{quote}
\item Olsson: \emph{Hom-stacks and Restrictions of Scalars}
\cite{olsson_homstacks}
\begin{quote}
$\mathcal{X}$ と $\mathcal{Y}$ を有限な対角射をもつ代数空間 $S$ 上
局所有限表示な Artin スタックとし，$\mathcal{X}$ は $S$ 上固有かつ
平坦で，$S$ 上 fppf 局所的に $\mathcal{X}$ が代数空間による有限，
有限表示，平坦な被覆を
許すとする（たとえば $\mathcal{X}$ が Deligne--Mumford スタックまたは
tame Artin スタックである場合）．このとき
$\Hom_S(\mathcal{X}, \mathcal{Y})$ は $S$ 上局所有限表示な
Artin スタックである．
\end{quote}
\item Alexeev and Knutson: \emph{Complete moduli spaces of branchvarieties}
(\cite{alexeev-knutson})
\begin{quote}
$\mathbf{P}^n$ の branchvariety を，\emph{被約}スキーム $X$ からの有限射
$X \rightarrow \mathbf{P}^n$ として定義する．Hilbert 多項式および
$i$ 次元成分の全次数を固定した branchvariety のモジュライ・スタックが，
有限安定化群をもつ固有 Artin スタックであることを証明する．
branchvariety のスタックを Hilbert スキーム，Chow スキーム，
安定写像のモジュライ空間と比較する．
\end{quote}
\item Lieblich: \emph{Remarks on the stack of coherent algebras}
\cite{lieblich_remarks}
\begin{quote}
この論文は，スキーム $Y$ 上の Alexeev--Knutson の branchvariety
スタックを一般化し，$Y$ の構造層上の代数のスタックとして構成する．
$\text{Quot}$ 空間と $\Hom$ 空間の存在証明が与えられる．
\end{quote}
\item Starr: \emph{Artin's axioms, composition, and moduli spaces}
\cite{starr_artin}
\begin{quote}
主結果の応用として，Vistoli の Hilbert スタック
\cite{vistoli_hilbert} と Alexeev--Knutson の branchvariety スタック
\cite{alexeev-knutson} の共通の一般化が与えられる．
$\mathcal{X}$ を優秀スキーム $S$ 上局所有限型で有限な対角射をもつ
代数スタックとする．このとき，$G$-豊富直線束 $L$ をもつ固有代数空間
$T$ からの射 $g: T \rightarrow \mathcal{X}$ をパラメータ化する
スタック $\mathcal{H}$ は，$S$ 上局所有限型な Artin スタックである．
\end{quote}
\item Lundkvist and Skjelnes:
\emph{Non-effective deformations of Grothendieck's Hilbert functor}
\cite{lundkvist-skjelnes}
\begin{quote}
非有効変形が存在するため，非分離スキームの Hilbert 関手が表現可能でない
ことを示す．
\end{quote}
\item Halpern-Leistner and Preygel:
\emph{Mapping stacks and categorical notions of properness}
\cite{HL-P}
\begin{quote}
始域と標的に関する，Olsson の論文の仮定より一般的，または少なくとも
異なるいくつかの仮定の下で，Hom スタックが代数的であることを証明する．
\end{quote}
\end{itemize}

\subsection{トーリック・スタック}
\label{guide-subsection-toric}

\noindent
トーリック多様体がスキーム論における例と反例を与えるのと同様に，
トーリック・スタックの幾何学と stacky fan の組合せ論との間には
対応関係がある．そのためトーリック・スタックは，豊富な例のクラスと，
予想を試す自然な場を与える．
\begin{itemize}
\item Borisov, Chen and Smith: \emph{The orbifold Chow ring of toric
Deligne-Mumford stacks} \cite{bcs}
\begin{quote}
この論文は，トーリック多様体に対する Cox の構成に着想を得て，
\emph{stacky fan} と呼ばれる組合せ論的対象に付随する明示的な
商スタックとして，滑らかなトーリック DM スタックを定義する．
\end{quote}
\item Iwanari: \emph{The category of toric stacks} \cite{iwanari_toric}
\begin{quote}
この論文は \emph{トーリック三つ組}を，稠密な像をもつ開埋め込み
$\mathbf{G}_m \hookrightarrow \mathcal{X}$（したがって
$\mathcal{X}$ は orbifold である）と作用
$\mathcal{X} \times \mathbf{G}_m \rightarrow \mathcal{X}$ を備えた
滑らかな Deligne--Mumford スタック $\mathcal{X}$ として定義する．
トーリック三つ組の 2-圏と stacky fan の 1-圏との間に同値があることを
示す．トーリック三つ組と \cite{bcs} における滑らかなトーリック
DM スタックの定義との関係も論じる．
\end{quote}
\item Iwanari: \emph{Integral Chow rings for toric stacks}
\cite{iwanari_chow}
\begin{quote}
トーリック多様体に対する Cox の $\Delta$-collection を
トーリック orbifold へ一般化する．
\end{quote}
\item Perroni: \emph{A note on toric Deligne-Mumford stacks}
\cite{perroni}
\begin{quote}
Cox の $\Delta$-collection と Iwanari の論文 \cite{iwanari_chow} を，
一般の滑らかなトーリック DM スタックへ一般化する．
\end{quote}
\item Fantechi, Mann, and Nironi: \emph{Smooth toric DM stacks}
\cite{fmn}
\begin{quote}
この論文は，滑らかなトーリック DM スタックを，DM トーラス
$\mathcal{T}$（すなわち $G$ が有限である $T \times BG$ と同型な
Picard スタック）の作用を備え，$\mathcal{T}$ と同型な稠密開軌道をもつ
滑らかな DM スタック $\mathcal{X}$ として定義する．
「ボトムアップ記述」を与え，滑らかなトーリック DM スタックと
stacky fan の間の同値を証明する．
\end{quote}
\item Geraschenko and Satriano: \emph{Toric Stacks I and II}
\cite{gs_toric1} および \cite{gs_toric2}
\begin{quote}
これらの論文はトーリック・スタックを，トーリック多様体をそのトーラスの
部分群で割ったスタック商として定義する．一般点で stacky な
トーリック・スタックを，トーリック・スタックのトーラス不変部分スタック
として定義する．この定義は，従来のトーリック・スタックの諸定義を包含し，
拡張する．第 1 論文は stacky fan の組合せ論と，対応するスタックの幾何学との
間の対応関係を展開する．また \cite{perroni} の解釈を一般化して，
滑らかなトーリック・スタックのモジュライ解釈を与える．
第 2 論文はトーリック・スタックの内在的特徴づけを証明する．
\end{quote}
\end{itemize}



\subsection{形式関数定理と Grothendieck の存在定理}
\label{guide-subsection-theorem-formal-functions}

\noindent
以下の論文は，形式関数定理 \cite[III.4.1.5]{EGA}
（固有射に対する Grothendieck の基本定理と呼ばれることもある）と，
Grothendieck の存在定理 \cite[III.5.1.4]{EGA} の一般化を与える．
\begin{itemize}
\item Knutson: \emph{Algebraic spaces} \cite[Chapter V]{Kn}
\begin{quote}
これらの定理を代数空間へ一般化する．
\end{quote}
\item Abramovich-Vistoli: \emph{Compactifying the space of stable maps}
\cite[A.1.1]{abramovich-vistoli}
\begin{quote}
これらの定理を tame Deligne--Mumford スタックへ一般化する．
\end{quote}
\item Olsson and Starr: \emph{Quot functors for Deligne-Mumford stacks}
\cite{olsson-starr}
\begin{quote}
これらの定理を分離的 Deligne--Mumford スタックへ一般化する．
\end{quote}
\item Olsson: \emph{On proper coverings of Artin stacks}
\cite{olsson_proper}
\begin{quote}
固有 Artin スタックへの一般化を与える．
\end{quote}
\item Conrad: \emph{Formal GAGA on Artin stacks} \cite{conrad_gaga}
\begin{quote}
固有 Artin スタックへの一般化を与え，形式 GAGA 定理を証明する．
\end{quote}
\item Olsson: \emph{Sheaves on Artin stacks} \cite{olsson_sheaves}
\begin{quote}
固有 Artin スタックへの一般化の別証明を与える．
\end{quote}
\end{itemize}


\subsection{スタック上の群作用}
\label{guide-subsection-group-actions}

\noindent
代数スタック上の群作用は自然に現れる．たとえば対称群 $S_n$ は
$\overline{\mathcal{M}}_{g, n}$ に作用し，群 $G$ がスキーム $X$ に
作用するとき，$\text{Aut}(X)$ における $G$ の正規化群は $[X/G]$ に
作用する．さらに，スタック上のトーラス作用は Gromov--Witten 理論に
しばしば現れる．
\begin{itemize}
\item Romagny: \emph{Group actions on stacks and applications}
\cite{romagny_actions}
\begin{quote}
この論文は，群が代数スタックに作用するとは何を意味するかを厳密化し，
代数スタック上の群スキームの作用について，不動点の存在と商の存在を
証明する．Romagny の以前のノート \cite{romagny_notes} も参照されたい．
\end{quote}
\end{itemize}

\subsection{直線束の根を取ること}
\label{guide-subsection-root-stacks}

\noindent
この有用な構成は Cadman と Abramovich--Graber--Vistoli によって
独立に発見された．有効 Cartier 因子 $D$ をもつスキーム $X$ が与えられた
とき，$r$ 次根スタックは $D$ に沿って $X$ 上分岐し，$D$ 上では
$\mu_r$ 安定化群をもち，$D$ の外ではスキームのように振る舞う
Artin スタックである．
\begin{itemize}
\item Charles Cadman
\emph{Using Stacks to Impose Tangency Conditions on Curves}
\cite{cadman}
\item Abramovich, Graber, Vistoli: \emph{Gromov-Witten theory for
Deligne-Mumford stacks} \cite{agv}
\end{itemize}

\subsection{その他の論文}
\label{guide-subsection-other}

\noindent
その他のさまざまな論文．
\begin{itemize}
\item Lieblich: \emph{Moduli of twisted sheaves} \cite{lieblich_twisted}
\begin{quote}
この論文はジェルブと twisted 層の要約を含む．
$\mathcal{X} \rightarrow X$ を $\mu_n$-ジェルブとし，
$X$ を滑らかで連結な幾何学的ファイバーをもつ射影的相対曲面とする．
このとき半安定な $\mathcal{X}$-twisted 層のスタックが，$S$ 上
局所有限表示な Artin スタックであることを示す．また，Artin スタック上の
層の随伴点と純粋性の理論も展開する．
\end{quote}
\item Lieblich, Osserman:
\emph{Functorial reconstruction theorem for stacks}
\cite{lieblich-osserman}
\begin{quote}
代数スタックが，それに付随する関手からいつ復元できるかについて，
驚くべき興味深い結果をいくつか証明する．
\end{quote}
\item David Rydh:
\emph{Noetherian approximation of algebraic spaces and stacks}
\cite{rydh_approx}
\begin{quote}
準有限な対角射をもつ任意の準コンパクト代数スタックが，
Noether スタックによって近似できることを示す．
Chevalley, Serre, Zariski, Chow の結果から Noether 性の仮定を除くことへの
応用がある．
\end{quote}
\end{itemize}

\section{他分野におけるスタック}
\label{guide-section-stacks-areas}

\begin{itemize}
\item Behrend and Noohi: \emph{Uniformization of Deligne-Mumford curves}
\cite{behrend-noohi}
\begin{quote}
位相的，解析的，代数的スタックの概観と比較を与える．
\end{quote}
\item Behrang Noohi: \emph{Foundations of topological stacks I} \cite{noohi}
\item David Metzler: \emph{Topological and smooth stacks} \cite{metzler}
\end{itemize}


\section{高次スタック}
\label{guide-section-higher-stacks}

\begin{itemize}
\item Lurie: \emph{Higher topos theory} \cite{lurie_topos}
\item Lurie: \emph{Derived Algebraic Geometry I - V}
\cite{dag1}, \cite{dag2}, \cite{dag3}, \cite{dag4}, \cite{dag5}
\item To\"en: \emph{Higher and derived stacks: a global overview}
\cite{toen_higher}
\item To\"en and Vezzosi: \emph{Homotopical algebraic geometry I, II}
\cite{hag1}, \cite{hag2}
\end{itemize}
